\chapter{多卷本路线}

本书不是要把周到清的两千多年压成一本厚得无法阅读的 PDF。每一卷都要有完整的事件账本、国家 dossier、跨国合卷、经济社会与文化思想章节、地图、史料争议和索引。

\begin{center}
\begin{tabularx}{\linewidth}{@{}l l X@{}}
\toprule
\textbf{卷} & \textbf{范围} & \textbf{核心问题} \\
\midrule
第一卷 & 西周至秦 & 封建网络如何裂解为领土国家，统一为何迅速崩解 \\
第二卷 & 汉至南北朝 & 帝国怎样调节秦制，又如何进入长期分裂 \\
第三卷 & 隋唐五代 & 再统一、世界帝国、藩镇、财政与社会转型 \\
第四卷 & 宋辽夏金元 & 多政权竞争、商业财政、征服与跨区域网络 \\
第五卷 & 明 & 皇权、赋役、白银、边疆与晚期危机 \\
第六卷 & 清至1912 & 征服王朝、人口与财政、近代冲击、革命与帝国终结 \\
\bottomrule
\end{tabularx}
\end{center}

第一卷的标准将决定后续所有卷：事件要足够密，国家线与跨国线要能互相跳转，制度得失要能落在具体人和具体年份上，经济社会与文化思想不能沦为附录。
